\documentclass[11pt]{article}
\usepackage[T1]{fontenc}
\usepackage[utf8]{inputenc}
\usepackage{lmodern,amsmath,amssymb,graphicx,geometry,microtype,setspace,booktabs,caption}
\usepackage[hidelinks]{hyperref}
\usepackage{array}
\geometry{letterpaper,margin=1in}
\setlength{\parindent}{0pt}
\setlength{\parskip}{0.55em}
\setlength{\emergencystretch}{2em}
\captionsetup{labelfont=bf,font=small}
\renewcommand{\refname}{References}
\begin{document}
\begin{center}{\Large\bfseries S2 Appendix: Data and sensitivity analyses}\\[0.5em]Yule Choi\end{center}
\onehalfspacing
\paragraph{Reading guide.} D1--D2 report digitization, fitting and sensitivity to the chosen response model. D3 examines prospective state measurements; D4 states which biochemical observations constrain the models; D5 specifies nuisance-aware assay precision; D6 documents the independent peptide data and normalization uncertainty. Mathematical mechanisms and proofs are in S1 Appendix. Earlier exploratory analyses and figures are retained in the repository archive instead of duplicating the main manuscript here. All synthetic calculations are distinguished from experimental measurements.
\section*{D1 Coordinate reconstruction and model comparison}

The canonical input is the author-supplied CSV digitized by eye from an enlarged view of Fig. 2B in the published first-layer study, as documented in the accompanying data provenance. It contains twelve rounded coordinate pairs. As an independent consistency check, vector-marker coordinates from page 19, Fig. 5b of the supplied manuscript were transformed using


\[
\log_{10}(l/{\rm mM})=\frac{X-436.74244}{489.11522-436.74244},\qquad
U=\frac{322.80995-Y}{322.80995-288.14189}. \tag{S21}
\]

Here X and Y are PDF points and Y is measured downward. The extracted coordinates agree with the supplied CSV to plotting precision. The unrounded extraction is retained as a consistency check, not as a more precise measurement. All reported numerical values agree at the displayed precision when fitted from either input. Neither source recovers experimental replicates or independently validates digitization against the original publication.


The free-exponent fit has SSE=0.01487834, h=2.017894, S=0.559542 mM, $U_{\min}$=0.453502, and $U_{\max}$=2.967708. One hundred random-start fits returned the same SSE to a spread below 5$\times$$10^{-16}$. This is a numerical convergence check, not a proof of a global optimum.


\begin{center}\begin{singlespace}\small
\begin{tabular}{lrrrr}\toprule

Model & SSE & k including variance & $\Delta$AIC & $\Delta$AICc \\
\midrule

h=1 & 0.527163 & 4 & 42.764 & 42.764 \\

h=2 & 0.014937 & 4 & 0 & 0 \\

h=3 & 0.119265 & 4 & 24.930 & 24.930 \\

h free & 0.014878 & 5 & 1.953 & 6.238 \\

\bottomrule\end{tabular}
\end{singlespace}\end{center}

The h=1 fit reaches the physical upper-plateau bound. Removing that bound gives $U_{\max}$=3.32303 and SSE=0.365830, reproducing the less penalized fit in the earlier analysis. Values above three cannot be interpreted as UMP occupancy per trimer. The displayed comparisons use the same physical bounds for every model and are descriptive model scores rather than tests of binding-site count.


The nominal profile set for h is approximately [1.851,2.203]. Under the specified 2,000 coordinate perturbations, its median is 2.025 and its central 95\% range is [1.750,2.421]. These intervals answer different questions: conditional fit curvature and sensitivity to an imposed coordinate-error distribution. No MCMC posterior or experimental confidence interval is claimed.


The regulatory-swing point estimate is 516.1 and the finite lower endpoint of the nominal profile set is 172.9. Its large-W tail remains inside the set. The coordinate perturbations give a 2.5th percentile of approximately 125 and an effectively unbounded upper range, with 519 fitted upper plateaus within $10^{-5}$ of three. The archive includes individual draws rather than summarizing this tail with a symmetric error bar. The reverse-activation ratio is treated as fixed; uncertainty in that assay is not included.


\section*{D2 Finite-switch fits and exact mean-matched ladders}

All new upstream comparisons use the twelve coordinates already documented in D1, including the author-supplied 0.02 mM first coordinate. The ordinary bounded Hill reference is $U(G)=b+(t-b)/(1+(G/S)^h)$. The bounds are $1.50001\le t\le2.999999$, $0.00001\le b\le1.49999$, $0.005\le S\le30$ mM, and $0.1\le h\le8$. The small numerical offsets keep inverse weights finite and exclude degenerate equal plateaus. They alter the original reference fit negligibly. The four optimized coordinates are $t,b,\log S,h$; inputs are not optimized as latent observations. The reference is fitted from twelve starts, each fixed-shape finite model from eight, and each profile point from two, using seeded starts and bounded least squares.

For the finite-switch models, $c=(1,3e^{-2J},3e^{-2J},1)$ and $m_J(u)=\sum_i i c_i u^i/(3\sum_i c_i u^i)$. Set $u_0=m_J^{-1}(t/3)$, $u_\infty=m_J^{-1}(b/3)$, and $u_m=m_J^{-1}((t+b)/6)$. Define
\[
z(G)=\frac{u_0-u_m}{u_m-u_\infty}(G/S)^h,\qquad
u(G)=u_\infty+\frac{u_0-u_\infty}{1+z(G)},\qquad U_J(G)=3m_J(u(G)). \tag{S36}
\]
This saturating forward/reverse weight is equivalent, by reciprocation and a scale change, to a basal saturating effective reverse/forward switch. It has the specified response endpoints and midpoint. The exponent is an effective regulation parameter, not ligand-binding stoichiometry. At $J=0$, Eq. (S36) equals the reference four-parameter Hill curve exactly. The fixed $J$ values are scenarios; they are not additional parameters covertly optimized in the three four-parameter fits. A separate 61-point profile over $[-1.5,1.5]$ is supplied as a sensitivity analysis and does not turn the chosen scenarios into independently selected models.

\begin{center}\begin{tabular}{rrrrr}\toprule $J$ & Fitted $h$ & RMSE & LOO RMSE & $\Delta$AICc \\\midrule
-0.7 & 3.115 & 0.0392 & 0.0634 & 2.56 \\
0.0 & 2.018 & 0.0352 & 0.0499 & 0.00 \\
0.7 & 1.414 & 0.0357 & 0.0503 & 0.31 \\
\bottomrule\end{tabular}\end{center}


Corrected AIC is $N\log(\mathrm{SSE}/N)+2k+2k(k+1)/(N-k-1)$, with $N=12$ and $k=5$, including an unknown common residual variance. Additive constants cancel between models. This working Gaussian residual comparison is not based on replicated assay errors. Leave-one-coordinate-out prediction refits all four parameters on eleven points using three starts; errors are in UMP groups per trimer. The 12 predictions for each model are supplied. The apparent similarity of these fits does not validate each molecular mechanism.

For the exact construction, any positive ladder satisfies
\[
\frac{d m_c}{d\log u}=\frac{\operatorname{Var}(i)}n>0,\quad
\lim_{u\to0}m_c=0,\quad\lim_{u\to\infty}m_c=1. \tag{S37}
\]
It therefore has a unique inverse, and $u_c(G)=m_c^{-1}(U(G)/3)$ gives the same full mean curve for each chosen ladder. For a decreasing mean, $u_c$ is decreasing and $f_c=1/u_c$ increasing. Positive birth--death rates with these adjacent ratios realize the distributions at each input; this does not prove a finite elementary ligand-binding implementation of the whole input function. Inversion is evaluated with a monotone interpolator and checked against independent bracketed roots. Exact mean matching is a structural construction, not a finite-dimensional fit assigned an AIC or evidence score.

The ladder parameter convention is $\log c_i=\log {n\choose i}+J[(i-n/2)^2-(n/2)^2]$. For $n=3$, this gives the coefficients above. Its sign describes endpoint or intermediate-state enrichment, not necessarily thermodynamic site cooperativity: effective coefficients can also contain state-specific enzyme recognition and turnover. Neither the uniqueness nor stability claims for the original single-effector architecture are assumed for these effective switch functions.

\section*{D3 State observables, prospective separation, and numerical checks}

For positive states in a fixed separable ladder,
\[
C_i=\frac{p_{i-1}p_{i+1}}{p_i^2}=\frac{c_{i-1}c_{i+1}}{c_i^2}. \tag{S40}
\]
These contrasts are unchanged by the input function or by the unavoidable rescaling $c_i\mapsto c_i a^i$, $u\mapsto u/a$. At a single input, state probabilities identify the ladder shape modulo that scale: $c_i c_0^{i-1}/c_1^i=p_i p_0^{i-1}/p_1^i$. At multiple inputs, changes in $C_i$ reject a fixed separable ladder for that observation model. State-dependent binding into experimentally counted complexes can also change total-state contrasts, so the test must be applied to the modeled observable.

For the symmetric three-site family, $3C_1=3C_2=e^{2J}$. In the refractory mixture, the first three states belong entirely to the accessible binomial population. Consequently $3C_1=1$, while direct substitution gives $3C_2=1+r/[(1-r)q^3]$. Increasing glutamine decreases $q$ and increases the latter contrast, although rare-state noise can make the ratio unsuitable near the plateau. Direct measurement of $p_3$ avoids division by vanishing intermediate probabilities.

For a practical separation calculation, let $D_{ab}(G)=\tfrac12\sum_{i=0}^3|p_i^{(a)}(G)-p_i^{(b)}(G)|$. We evaluate all six model pairs on 401 logarithmically spaced inputs from 0.02 to 10 mM and maximize $\min_{a<b}D_{ab}(G)$. The fixed candidates are the three exact mean-matched ladders and the refractory mixture. The selected input is 0.746758 mM and the minimum distance 0.159647. This is a maximin separation of latent trimer fractions, not a statistical power calculation or an optimum over all plausible mechanisms. No detection limit, band-specific bias, or instrument error covariance has been estimated.

At this input and at 10 mM, 300 shared coordinate draws generate pointwise prediction percentiles for every state. At 10 mM, the independent and refractory $p_3$ medians are 0.00362 and 0.15115, with central ranges [0.00178,0.00670] and [0.11660,0.18711]. These ranges omit uncertainty in the fixed ladder and population assumptions. A measurement near the transition and another at high input address different ambiguities: distribution shape and the origin of the response floor, respectively. The complete state probabilities and draw-level results accompany S1 Data.

Verification uses independent mathematical routes where available. Population equivalence is checked for 1,200 random cases over $n=1,\ldots,12$, with maximum relative identity error below $7\times10^{-13}$. Interpolated inverse weights agree with bracketed roots within $2.1\times10^{-9}$ in log weight over the verification grid. Finite differences of the mean agree with the variance derivative within $5.2\times10^{-11}$; the adjacent-state contrasts agree within $4.5\times10^{-15}$. An independent analytic gain calibration followed by bracketed threshold roots gives baseline coefficients 3.55184499, 3.77942277, and 4.91450595. A 20,001-point numerical calculation differs by at most $1.4\times10^{-5}$; population-equivalent exact coefficients differ by less than $5\times10^{-15}$. The generated verification file records these tolerances.

\clearpage\renewcommand{\thefigure}{S\arabic{figure}}\setcounter{figure}{0}

\section*{D4 What biochemical measurements constrain, and what remains unidentified}

Jiang, Peliska and Ninfa (1998), Table 2, reports Mg-dependent UT glutamine inhibition at 70--80 $\mu$M and UR activation at 80 $\mu$M. The Mg-dependent UR specificity ratio calculated from the rounded $k_{\rm cat}$ and $K_m$ entries is $(6.5/0.82)/(2.7/2.3)=6.752484$. The older manuscript's rounded specificity entries 7.93 and 1.17 give 6.778 instead; this small arithmetic difference is retained explicitly rather than treated as a new measurement. The Mn-dependent entries are UT inhibition 150 $\mu$M and UR activation 0.70 mM, with the footnoted cofactor restrictions. Correcting transposed Mn metadata and the compensating plotting ratio preserves the numerical ceiling; S1 Data records both corrections.

These measurements are apparent activity constants and catalytic summaries under their stated substrate/cofactor conditions. They are not measurements of the microscopic sequential ligand-binding constants in A5. In the independent-binding fit, transporting only the UR specificity fold gives an RMSE of 0.0366, but UT and UR half-ranges of approximately 0.022 and 0.126 mM. In the fixed-$C$ alternative, transporting the UR fold and 0.080 mM half-range gives an RMSE of 0.0378. A stricter test also transports the unliganded specificity ratio $(137/3)/(2.7/2.3)=38.9012$, fixing $q_0=0.974938$. Refitting the remaining parameter in that fixed-$C$ topology gives an RMSE of 0.540. Thus the partial success must not be presented as agreement with the full enzyme-kinetic dataset. Nor does failure of this cross-assay transfer rule out every finite model.

Jiang, Mayo and Ninfa (2007), Tables 2--3, reports PII apparent activation constants from 0.015 to 2.5 $\mu$M in the selected rows without PII-UMP, as glutamine and other conditions change. Rows including PII-UMP reach approximately 6 and greater than 40 $\mu$M. Reported PII-UMP activation constants include 0.28, 0.35, and 0.85 $\mu$M under different conditions. These motivate a broad scenario envelope; pooling them does not produce an experimentally justified prior or confidence interval for $K_P$ or $K_Q$. Source-row identifiers and conditions are retained in S1 Data. Independent receptor sites are a candidate simplification.

Write $P=T p_0$ and $Q=T\theta$. For a finite receptor with separate regulatory sites, the four aggregated occupancy classes have weights $(1,P/K_P,Q/K_Q,PQ/(K_PK_Q))$. The reverse input $Q$ is shorthand for three competing ligands $S_1,S_2,S_3$ with dissociation constants $3K_Q/i$ for state $i$ and equal effects after binding. The underlying receptor thus has eight microstates (forward site empty/bound, reverse site empty/bound to one of three states). This is a count-weighted recognition assumption. Adding identical glutamine regulation on both sides of the comparison preserves the following identity as long as the normalized input weights are preserved; it does not require PII-independent AT activity to be zero.

The population construction gives $P_{\rm R}=cP_{\rm H}$ with $c=(1-r)^{-2}$ and leaves $Q$ unchanged. Setting $K_{P,\rm R}=cK_{P,\rm H}$ preserves every occupancy class. More generally, multiplying every dissociation constant for unmodified PII by $c$ preserves the corresponding binding weights, including condition-specific ones. In the independent-binding fitted pair, $c$ is approximately 1.37. This is one shared transformation over all glutamine and target levels, not a separate calibration at each data point. It establishes an identifiable combination of population composition and affinity under the chosen observation model. It does not demonstrate that either parameter set is biologically correct.

If $K_P$ is known and held fixed, a separate-site saturated ratio $v=k[P/(K_P+P)]/[Q/(K_Q+Q)]$ is generally not preserved by a constant gain. The supplied 663 scenarios use three PII totals, 17 $K_P$ values, and 13 $K_Q$ values. At each setting, the homogeneous gain makes $Y=v/(1+v)=1/2$ at 0.53 mM and the refractory gain matches that response. The output is the maximum absolute GS-fraction difference across 401 inputs from 0.02 to 10 mM; it is not a response-range coefficient or a probability of experimental discrimination. Allowing the shared affinity transformation restores the full curve to numerical precision. State-specific reverse recognition, ligand/target sequestration, or independently fixed affinities can invalidate this result. No closed-total GlnE/PII conservation fit is claimed.

Discriminating measurements are matched-condition UT/UR initial rates, high-input PII state fractions, and state-resolved GlnE binding. These remain prospective tests.


\clearpage

\paragraph{Assay transfer audit and source-anchored benchmark.}
The 1998 Mg-UT assay uses 100 mM Tris, pH 7.5, 100 mM KCl, 25 mM MgCl$_2$, 0.5 mM ATP and a default 33 $\mu$M 2-OG; Mg-UR uses 50 mM MgCl$_2$, 0.5 mM ATP and default 0.1 mM 2-OG. Both are at $30^\circ$C, but their substrate and effector conditions differ. The 2011 coupled assay uses 50 mM Tris, 10 mM MgCl$_2$, 100 mM KCl, 1 mM ATP, 0.5 mM UTP, 1 mM phosphate and 0.1 mM 2-OG at the same pH and temperature. Its PII/GlnD pairs are $(0.5,0.05),(5,0.5),(36,0.8)\,\mu$M, with GlnE 0.1 $\mu$M. It verifies stationary values using multiple time points. These differences prevent interpreting every apparent specificity ratio as a microscopic coefficient already validated in the coupled system. The proposed concentration example explicitly counts GS in dodecamers; monomer/dodecamer conversions must be checked when transferring other assay quantities.

The preceding RMSE 0.540 construction fixes $C=100$, $q_1=q_0$, two endpoint specificity ratios and the UR half-range at 0.080 mM (A5). Its state capacities are $(1,1,d_2)$, so they are not all equal. The new fit retains $a_0/b_0=38.9012346$ and $b_2/b_0=6.7524842$, sets $b_0=1$, and imposes monotone $a_0\geq a_1\geq a_2$ and $b_0\leq b_1\leq b_2$. It releases the fixed-$C$, intermediate-state and 0.080 mM half-range constraints together; the RMSE improvement cannot isolate their individual effects. Fifty starts minimize unweighted residuals at the twelve digitized points, with $\log K_1,\log K_2\in[-12,5]$ (mM) and bounded logistic coordinates for monotone capacities. RMSE is 0.0409346, with $K_1=148.413$ mM at the upper bound and $K_2=0.000443$ mM. The UR capacity $R(G)=\sum_j b_jw_j/\sum_j w_j$ reaches $(b_0+b_2)/2$ at 0.256624 mM, confirming that the transported half-range is not retained. This boundary optimum is not evidence for a plausible affinity. The exact transformation remains admissible for
\[
 \max(a_2/a_1,b_0/b_1)\leq\lambda\leq\min(a_0/a_1,b_2/b_1).
\]
It preserves the transported endpoint ratios as well as the response, while changing the intermediate state. A profile fixing $K_1$ uses twelve starts per value and the same monotonic constraints. At $K_1=0.3,1,10,100$ mM, RMSE is $0.0493,0.0430,0.0411,0.0409$ and maximum singly liganded contrast at $G=\sqrt{K_1K_2}$ is $0.242,0.0954,0.0108,0.00109$. This contrast is the width over the admissible transformed family, not a confidence interval or a claim that every pair is distinguishable. The profile is a sensitivity analysis, not a global likelihood confidence calculation. Source ratios transferred across assays remain stress assumptions.

The GlnE benchmark is more directly tied to a reported fit. Jiang et al. (2007) supplement \href{https://doi.org/10.1021/bi0620510.s001}{doi:10.1021/bi0620510.s001}, Fig S10 (PDF p.16), uses GlnE 0.04 $\mu$M, GS 2.5 $\mu$M, ATP 0.5 mM and 2-OG 0.05 mM. It reports $K_G=15.6$ mM, $\alpha_1=0.17$ and EP/EG/EGP catalytic coefficients $7.5/175/1315$, giving $\beta=23.3333$ and $\beta'=175.3333$. The main article states $K_P=0.18\,\mu$M and a sixfold reduction relative to the direct apparent scale, whereas the S10 legend prints 0.18 mM. We use the main-text unit and explicitly calculate the literal supplementary alternative. Neither is an independently measured microscopic uncertainty distribution. The $K_G\in[7.8,15.6]$ mM benchmark spans the main article's approximately twofold direct-scale comparison and joint fit, not a measured confidence region. The 2007 s002 supplementary file has a different figure sequence and is not substituted for s001.

The main concentration example combines these source-fit coefficients with reconstitution-motivated protein scales as a declared prospective scenario. It does not claim that the assays constitute one matched dataset. A concurrently measured PII count upper bound and an independently supported coefficient envelope are the operational requirements for applying A11. In particular, the rounded 0.65 coordinate at 2 mM does not establish $M\leq0.67$ with any confidence level. The older illustrative coefficients $\beta=1,\beta'=10$ remain only in the original stress sweeps documented in A10.

\section*{D5 Binding observables and assay performance}

\textbf{GlnD occupancy.} In a target-free equilibrium binding assay,
\[
 o_1(z,\lambda)=\frac{2z/\lambda}{1+2z/\lambda+z^2},\qquad z=G/K.
\]
At nominal $G=K_0$, allow $K/K_0\in[1/2,2]$ separately under each hypothesis. For $\lambda=1$, the fraction lies in $[4/9,1/2]$; for $\lambda=4$, it lies in $[1/6,1/5]$. The minimum gap is $4/9-1/5=11/45$. This interval separation profiles the nuisance ranges independently, rather than imposing the same unknown $K$ on the two candidates. The catalytic scale and substrate affinities do not enter this target-free assay. The nominal concentration is approximately 0.0520 mM from the reference fit; the factor-two range is a design scenario, not an inferred confidence interval.

For a general alternative $\lambda\geq1$, its largest single-state occupancy over the same range is $1/(1+\lambda)$. The worst-case gap from the reference is $4/9-1/(1+\lambda)$, positive only when $\lambda>5/4$. When ranges overlap, replication of this scalar measurement alone cannot give uniform discrimination over the nuisance classes. Multiple glutamine levels or a tighter independently supported $K$ range are then required.

\textbf{GlnE binding and depletion.} Set PII-UMP to zero and keep glutamine buffered. The three PII-bound states reduce to the unmodified-PII and glutamine--PII complexes, and the PII-bound enzyme fraction is $P_f/(K_{\rm eff}+P_f)$, where
\[
 K_{\rm eff}=K_P\frac{1+G/K_G}{1+G/(\alpha_1K_G)}.
\]
If $P_T=P_f+E_T P_f/(K_{\rm eff}+P_f)$, half occupancy implies $P_T^{1/2}=K_{\rm eff}+E_T/2$. Thus a measured enzyme total corrects the nominal halfpoint before normalization across glutamine conditions. This correction assumes one PII binding unit per enzyme in the stated topology and does not correct glutamine depletion; the latter must be prevented or modelled separately.

The normalized ratio $r(G)=K_{\rm eff}(G)/K_{\rm eff}(0)$ eliminates $K_P$. For $\alpha_1\ne1$,
\[
 \frac{G}{1-r(G)}=\frac{\alpha_1K_G}{1-\alpha_1}+\frac{G}{1-\alpha_1}.
\]
Two distinct nonzero glutamine levels plus the zero-glutamine baseline determine the line's slope $m$ and intercept $b$, hence $\alpha_1=1-1/m$ and $K_G=b(1-\alpha_1)/\alpha_1$. At $\alpha_1=1$, $r=1$ and $K_G$ is not identifiable from this observation. This is a noise-free identifiability statement; noisy titration data should be fitted with the original binding model because transforming ratios changes their errors and shares the baseline error. Numerical inversion at 1 and 10 mM checks the formula across the declared parameter range.

For prospective discrimination at 0 and 10 mM, use two disjoint effect-size classes $\alpha_1\in[0.136,0.2125]$ and $[0.544,0.850]$. In both, allow $K_G\in[7.8,31.2]$ mM independently. This factor-two sensitivity range around 15.6 mM belongs to the earlier known-variance illustration. The observation benchmark below instead uses $[7.8,15.6]$ mM to span the direct-scale comparison and source fit (D4); its Monte Carlo results do not extend to the wider range. Monotonicity of the ratio gives log-ratio intervals $[-1.51931,-0.64158]$ and $[-0.38589,-0.04194]$. The minimum gap is 0.25570. These classes specify the alternatives being tested; neither is a confidence region from the published experiments. If each nominal $\alpha_1$ instead spans a factor of two, the classes touch and there is no strictly positive worst-case guarantee. The three-titration design then provides an alternative to an unjustified binary claim.

\textbf{Earlier analytic budget: known variance.} Let each occupancy estimate have independent normal error with known SD $\sigma$, or let each independent halfpoint estimate have normal log error with known SD $\sigma$. Background-corrected fraction estimates are not clipped in this error model. Precision must be established in pilot measurements; this analytic special case does not include uncertain error variance or condition-dependent calibration bias; the end-to-end benchmark below addresses both. A single GlnE halfpoint estimate is obtained from an entire titration, not one concentration point. For $n$ independent estimates per condition, the occupancy mean has variance $\sigma^2/n$, while the difference of two mean log halfpoints has variance $2\sigma^2/n$.

Use the least-favourable null boundary and the prespecified sign of the alternative. A one-sided size-0.05 normal test has worst-case power
\[
 \Phi\!\left(\frac{\sqrt n\,\Delta}{\sqrt v\,\sigma}-z_{0.95}\right),
 \qquad n\geq\frac{v\sigma^2(z_{0.95}+z_{0.80})^2}{\Delta^2},
\]
rounded upward. Applied to total-readout separation $\Delta\leq\epsilon$, it gives an optimistic necessary sample size $n\geq\sigma^2(z_{0.95}+z_{0.80})^2/\epsilon^2$ for one normally distributed fraction observable and two specified alternatives. With $\epsilon=0.01,\sigma=0.05$, this rounds up to 155; smaller actual separation or nuisance uncertainty can only worsen this bound. For the binding designs use $v=1$ for occupancy and $v=2$ for paired log halfpoints. Here $\Delta$ is the independently profiled interval gap, not a distance between two selected curves. At SD 0.20, five occupancy estimates yield worst-case power 0.862, and eight halfpoint estimates per condition yield 0.819. At occupancy SD 0.10, two estimates yield 0.965. The selected tests are independently verified by 100,000 synthetic null and alternative draws. These are prospective assay-error calculations, not experimentally established budgets. If independent full titrations are infeasible, a hierarchical joint fit needs a separate power analysis.

\textbf{Conditional kinetic comparison.} The capacity multiplier is $R_\lambda(z)=(1+2z+z^2)/(1+2z/\lambda+z^2)$. An unknown speed absorbs a single-condition amplitude. Paired capacities cancel a shared speed, but reciprocal doses are uninformative because $R_\lambda(z)=R_\lambda(1/z)$. Searching 101 doses over $0.02K_0$--$50K_0$ selects approximately $0.020K_0$ and $0.855K_0$, with minimum log-multiplier contrast 0.363 over factor-two $K$ variation. This contrast requires independently calibrated baseline capacities and $K$; independently refitting the baseline shape can restore overlap. The data label its sample sizes conditional, and Fig 4 excludes them from the profiled guarantees.



\paragraph{Observation-to-decision benchmark.}
The proposed observable is a solution signal proportional to the fraction of GlnE bound by PII, with PII-UMP absent. Free and saturating-ligand controls, label perturbation checks and glutamine-dependent signal controls would be required experimentally. The simulation assumes this readout, rather than asserting measured instrument performance. For free halfpoint $H$ and total enzyme $E$, the exact one-class bound fraction at total PII $P$ is
\[
 f(P;H,E)=\frac{2P}{P+H+E+\sqrt{(P+H+E)^2-4EP}}.
\]
It follows from $P=P_f+Ef$ and $f=P_f/(H+P_f)$; at half occupancy $P=H+E/2$. Halfpoints follow $H(G)=0.18(1+G/K_G)/(1+G/(\alpha_1K_G))\,\mu$M.

Each independent experiment contains $G=0,15.6,31.2$ mM and total PII doses 0, 0.005, 0.01, 0.02, 0.05, 0.1, 0.2, 0.5, 1, 2 and 5 $\mu$M, with two technical readings per dose. True $E=0.04\,\mu$M uses the source enzyme scale. In nonideal scenarios, one shared PII-stock multiplier has log SD 0.10 and the estimated enzyme total has log SD 0.20. One shared background has SD 0.01 and one shared signal gain has log SD 0.03 in every scenario. Independent observation noise has SD $\sqrt{0.025^2+(0.02f)^2}$; signals are not clipped. ``Ideal'' means exact dose and enzyme calibration, not noiseless observations.

The dose-bias scenario additionally multiplies the actual nonzero-glutamine PII doses by $\exp(-0.10)$, a systematic error that remains on replication. In the corrected scenario, independent concentration-standard readings estimate this factor for each nonzero-glutamine solution in each experiment, with log SD 0.02; these measured doses are used in fitting. The shared stock error and enzyme-total error remain. These magnitudes are declared sensitivity settings, not estimates from the source papers.

Three log halfpoints and a shared background/gain are fitted jointly by bounded nonlinear least squares, accounting for finite depletion with the estimated enzyme total. Starting halfpoints are $0.2,0.1,0.1\,\mu$M; allowable halfpoints are $[10^{-4},10]\,\mu$M, background $[-0.2,0.2]$ and gain $[0.5,1.5]$. Fits stop with $10^{-8}$ convergence tolerances or 150 evaluations. A failed optimizer or active parameter bound invalidates the entire trial and is counted as a nondiscovery. Every scheduled fit in the reported main and local-effect grids is valid.

For experiment $j$, define $R_{j1}=\log(\widehat H_{j1}/\widehat H_{j0})$, $R_{j2}=\log(\widehat H_{j2}/\widehat H_{j0})$ and $Z_j=(R_{j1}+R_{j2})/2$. Eight independent experiments give $\overline Z$ and sample SD $s_Z$. Under the reference $\alpha_1\leq0.17$, with $K_G\in[7.8,15.6]$ mM, the largest mean score in the noiseless one-class model is
\[
 z_0=\max_{K_G}\frac12\sum_{k=1}^2\log\frac{1+G_k/K_G}{1+G_k/(0.17K_G)}=-1.34194257.
\]
Both ratios increase with $K_G$ for $\alpha_1<1$, so the maximum is at 15.6 mM. Reject when $(\overline Z-z_0)/(s_Z/\sqrt8)>t_{0.95,7}$. This approximate test is evaluated empirically through the complete nonlinear estimation pipeline; a Student distribution is not asserted exactly for fitted halfpoints. The shared baseline covariance is retained because $s_Z^2=(s_{R1}^2+s_{R2}^2+2s_{R1,R2})/4$. The supplied comparison omits the covariance term to document the effect of a naive independence assumption. At the noise-only null boundary the mean within-trial ratio correlation is 0.652.

Five scenarios, three true $K_G$ values $7.8,11.7,15.6$ mM and $\alpha_1=0.17,0.25,0.34,0.50,0.68$ give 75 configurations, each with 100 Monte Carlo trials. Twelve additional corrected-dose configurations use $\alpha_1=0.18,0.19,0.20,0.22$ and the same three $K_G$ values. Seeds are 9142026 and 9142027. All 8,700 trial summaries and aggregate counts are supplied; each trial contains eight complete experiments, or 528 technical readings. Wilson intervals quantify Monte Carlo uncertainty only.

At $\alpha_1=0.17,K_G=15.6$ mM, rejection counts per 100 trials are 6 (ideal calibration), 3 (stock and enzyme errors), 96 (systematic dose bias), 0 (independently corrected bias), and 100 (the misspecified mixture below). Corresponding 95\% Wilson intervals are $[0.028,0.125]$, $[0.010,0.085]$, $[0.902,0.984]$, $[0,0.037]$ and $[0.963,1]$. These counts do not prove a universal false-positive guarantee. At corrected $\alpha_1=0.20$, rejection is 0, 0.46 and 0.98 as true $K_G$ increases across the grid; at $\alpha_1=0.25$ it is 1 in all three tested settings. Thus nuisance variation still controls small-effect power. Interpolating plotted points does not establish performance between tested settings.

The misspecified generator has two affinity classes with halfpoints $H$ and $5H$. Their high-affinity fraction changes from 0.9 at zero glutamine to 0.6 at nonzero glutamine. At each dose, the free ligand solves
\[
 P=P_f+E\left[w\frac{P_f}{H+P_f}+(1-w)\frac{P_f}{5H+P_f}\right].
\]
The data are nevertheless fitted with one class. At the reference boundary this gives 100/100 rejections, while mean fit RMSE is 0.0279, close to 0.0267--0.0271 for the four other scenarios. The changed class weights are an unmodeled regulatory mechanism; a single-class rejection cannot identify a specific change in $\alpha_1$. This deliberately tests the observation-model boundary rather than claiming robustness to all heterogeneity.

Finally, singly liganded GlnD is a more demanding prospective observable. It requires discrimination of zero-, one- and two-glutamine enzyme states. A native intact-protein MS pilot would need apo/ligand assignments, titration and ACT controls, and checks for ligand loss or adducts. GlnD state resolution remains unvalidated. The original 0.244 gap and five-estimate budget therefore remain illustrative rather than the central assay recommendation.

\section*{D6 Independent peptide data and normalization uncertainty}
The unchanged source workbook and model accompany the Figshare dataset of Gosztolai et al. (2017), doi:10.6084/m9.figshare.4880003, under CC BY 4.0. The article is in Biophysical Journal, doi:10.1016/j.bpj.2017.04.012. Extraction retains source cells, means, standard errors and sample counts in \texttt{observations.csv}. The source model converts protein copies to micromolar with factor 0.00161; the workbook label is retained and the conversion is documented rather than silently interpreted as molar. Peptide measurements quantify modified and unmodified subunits, not intact oligomer distributions.

For modified signal $U$, unmodified signal $V$ and a separately measured total $T$, paired normalization is $U/(U+V)$ whereas separate-total normalization is $U/T$. These are different observation models. Five separate-total point ratios exceed one; they are recorded as infeasible fractions rather than treated as valid occupancies. Unknown replicate covariance prevents reconstructing a joint likelihood from the published summary statistics. First-order SE bounds in Fig 5A and interval bounds in Fig 5D have different meanings.

Figure 5D uses the original component means, SE and sample counts. At each time point, each of three component means receives a two-sided working $t$ interval with tail probability $0.05/6$ and its own $n-1$ degrees of freedom; intervals are truncated to nonnegative values. Bonferroni adjustment permits dependence between the component estimates, conditional on valid marginal intervals. It does not establish Gaussian sampling, calibration validity, or coverage across all time points. Ratio extrema are evaluated over this box. At WT $t=-15$ min, paired normalization yields $q=0.9675734$ and lower bound $0.8384118$, hence ceilings $0.0324266$ and $0.1615882$. Separate-total normalization has lower bound $0.3319311$, hence ceiling $0.6680689$. The large difference is retained rather than choosing the narrower observation model as if independently validated. Linear programs verify all reported count-state extrema (624 optimizations for 156 records).

Neither normalization establishes a stationary free-pool response or the transferability of a no-binding calibration. The observed data therefore support conditional ceilings and prospective assay feasibility. The curve-calibration construction in A14 requires additional matched controls; its random-case numerical checks are not biological validation. A direct AmtB--GlnK binding endpoint, held out from calibration, would test the inferred interval. The repository document \texttt{BIOLOGICAL\_VALIDATION\_PLAN.md} separates these measurement roles.
\section*{D7 Reciprocal regulation and plateau identities}

For f(l)=C l(1+$A_{UT}$ l)/(1+$A_{UR}$ l), direct logarithmic differentiation gives $\eta=1+A_{UT}l/(1+A_{UT}l)-A_{UR}l/(1+A_{UR}l)$. If $A_{UT}$>$A_{UR}$, differentiating once more locates the maximum at l=$1/\sqrt{A_{UT}A_{UR}}$, where


\[
\eta_{\max}=\frac{2\sqrt{A_{UT}}}{\sqrt{A_{UT}}+\sqrt{A_{UR}}}. \tag{S14}
\]

If the affinities are equal, $\eta$=1. If $A_{UT}$<$A_{UR}$, the stationary point is a minimum and the supremum is one at the limiting inputs. The $\mathrm{Mg}^{2+}$ and $\mathrm{Mn}^{2+}$ values in S2 Fig B use inverse reported kinetic concentrations as model affinities: 80/70 and 700/150 for the two ratios. This is an explicit model identification, not an independent equilibrium-binding measurement.


On an independent ladder, take f=w+(v-w)$l^2$/($K_L^2$+$l^2$), with v>w>0. Substitution into $\theta$=K/(K+f) gives $\theta_0$=K/(K+w), $\theta_\infty$=K/(K+v), and


\[
\theta=\theta_\infty+\frac{\theta_0-\theta_\infty}{1+l^2/S^2},\qquad
S^2=K_L^2\frac{K+w}{K+v}. \tag{S15}
\]

Taking the ratio of the two plateau odds gives v/w. If v/w is the change in reverse-to-forward capacity, it equals reverse activation times forward inhibition. This identity requires a regulatory interpretation of both plateaus and the stated ladder. It does not hold as a mechanistic interpretation when the measured floor is instead a refractory fraction of target.


For a general fitted Hill curve with plateaus a and b, b>a, exponent h, and x=$(l/S)^h$, $\theta$=a+(b-a)/(1+x). Its physical local coefficient is


\[
n_H^{\rm loc}(l)=\frac{h(b-a)x}{(1+x)^2\theta(1-\theta)}. \tag{S16}
\]

At the fitted midpoint, this gives 1.725 for the reconstructed curve, whereas its range coefficient is 2.018. S2 Fig C displays this distinction directly. The fitted exponent coincides with a ligand count only in a separately specified regulatory limit.



\end{document}
